\subsection{关联张量丛与高阶协变导数}

局部标架表示规定场的纤维类型；联络将这一表示变成跨点的微分运算。

前面的群表示分解描述的是全局等距变换。若要研究场本身的自旋，还必须回到每一点的切空间，弄清张量怎样在不同点之间进行比较。

令 $\operatorname{Fr}_{SO}(M)$ 为定向正交标架主丛，结构群 $SO(d)$。给定其有限维表示 $\rho:SO(d)\to GL(V)$，定义关联向量丛
\begin{equation}\label{eq:ghs-associated-vector-bundle}
 E_\rho=\operatorname{Fr}_{SO}(M)\times_\rho V.
\end{equation}
它的截面 $\Gamma(E_\rho)$ 是空间上的场。对于整数自旋 $S$ 的对称无迹场，可取
\[
 E_S=\operatorname{Sym}^S_0(T^*M),\qquad
 \mathcal J^{(S)}\in\Gamma(E_S).
\]
当 $d>2$ 时，纤维维数为 $\binom{d+S-1}{S}-\binom{d+S-3}{S-2}$，第二项在 $S<2$ 时取零；这是对称张量空间减去迹像的维数，《群论》的调和张量章节给出完整分解，且与位置无关。物理自由度是否等于这一纤维维数，仍取决于运动方程、约束和规范对称性。

自旋量场需要更严格的前提：存在自旋结构时，正交标架丛才能提升到 $\operatorname{Spin}(d)$ 主丛。并非每个定向流形都有自旋结构；其障碍由第二 Stiefel--Whitney 类 $w_2(M)$ 表征。因此“任意空间上的任意自旋场”不能在不区分整数自旋与半整数自旋的情况下成立。

Levi--Civita 联络自然延拓到所有张量丛。对向量场 $V^c$，我们规定曲率号约定
\begin{equation}\label{eq:ghs-curvature-commutator}
 [\nabla_a,\nabla_b]V^c=R^c{}_{dab}V^d.
\end{equation}
对协向量使用度规相容性与 Leibniz 规则可得
$[\nabla_a,\nabla_b]\omega_c=-R^d{}_{cab}\omega_d$。推广到秩 $S$ 协变张量时，右端是 $S$ 个曲率项之和；每个指标单独贡献一次曲率作用。

\begin{proposition}[标量二阶导数与三阶导数]
对任意标量 $f$，$[\nabla_a,\nabla_b]f=0$；但在约定 \cref{eq:ghs-curvature-commutator} 下
\begin{equation}\label{eq:ghs-scalar-third-derivative}
 (\nabla_a\nabla_b-\nabla_b\nabla_a)(\nabla_cf)
 =-R^d{}_{cab}\nabla_df.
\end{equation}
\end{proposition}
\begin{proof}
标量的一阶协变导数就是偏导，故前两次作用交换。其梯度却是协向量，于是第三次求导时按照协向量的曲率对易关系获得右端。
\end{proof}

\begin{figure}[htbp]
\centering
% Local TikZ geometry; styles inherited exclusively from the shared preamble.
\begin{tikzfig}[scale=1.05]
\coordinate (O) at (0,0);\coordinate (A) at (2.65,.35);\coordinate (B) at (3.4,1.85);\coordinate (C) at (.75,1.5);
\fill[gray!7] (O)--(A)--(B)--(C)--cycle;
\draw[thick,lecturearrow] (O)--(A) node[midway,below] {$u$};
\draw[thick,lecturearrow] (A)--(B) node[midway,right] {$v$};
\draw[thick,lecturearrow] (B)--(C);
\draw[thick,lecturearrow] (C)--(O);
\fill (O) circle (1.6pt);
\draw[blue!65!black,lecturearrow] (O)--(.13,.9) node[left] {$V$};
\draw[red!65!black,lecturearrow] (O)--(.4,.81) node[right] {$V+\delta V$};
\node[lecturelabel] at (1.6,-.8) {$\delta V\propto R(u,v)V\;\cdot\;\mathrm{Area}$};
\end{tikzfig}

\caption{小闭合回路上的平行移动示意。沿两个方向交换传输顺序后，向量的差在最低非平凡阶由曲率 $R(u,v)$ 乘以有向面积控制。}\label{fig:ghs-curvature-holonomy}
\end{figure}

\cref{fig:ghs-curvature-holonomy} 中画出的两条末态箭头不是额外的自由度，而是同一向量在不同微小路径下平行移动的比较；当曲率为零时，它们在最低阶重合。

这说明“把平直空间公式中的所有 $\partial$ 换成 $\nabla$”不保证新公式守恒。高自旋流经常包含三阶及更多阶导数，曲率对易项在散度计算中不可避免。在一般曲率背景上，须同时考虑曲率修正、背景方程及可能的几何障碍。
